% Options for packages loaded elsewhere
\PassOptionsToPackage{unicode}{hyperref}
\PassOptionsToPackage{hyphens}{url}
\documentclass[
  11pt,
]{article}
\usepackage{xcolor}
\usepackage[margin=2,8cm]{geometry}
\usepackage{amsmath,amssymb}
\setcounter{secnumdepth}{5}
\usepackage{iftex}
\ifPDFTeX
  \usepackage[T1]{fontenc}
  \usepackage[utf8]{inputenc}
  \usepackage{textcomp} % provide euro and other symbols
\else % if luatex or xetex
  \usepackage{unicode-math} % this also loads fontspec
  \defaultfontfeatures{Scale=MatchLowercase}
  \defaultfontfeatures[\rmfamily]{Ligatures=TeX,Scale=1}
\fi
\usepackage{lmodern}
\ifPDFTeX\else
  % xetex/luatex font selection
\fi
% Use upquote if available, for straight quotes in verbatim environments
\IfFileExists{upquote.sty}{\usepackage{upquote}}{}
\IfFileExists{microtype.sty}{% use microtype if available
  \usepackage[]{microtype}
  \UseMicrotypeSet[protrusion]{basicmath} % disable protrusion for tt fonts
}{}
\makeatletter
\@ifundefined{KOMAClassName}{% if non-KOMA class
  \IfFileExists{parskip.sty}{%
    \usepackage{parskip}
  }{% else
    \setlength{\parindent}{0pt}
    \setlength{\parskip}{6pt plus 2pt minus 1pt}}
}{% if KOMA class
  \KOMAoptions{parskip=half}}
\makeatother
\usepackage{graphicx}
\makeatletter
\newsavebox\pandoc@box
\newcommand*\pandocbounded[1]{% scales image to fit in text height/width
  \sbox\pandoc@box{#1}%
  \Gscale@div\@tempa{\textheight}{\dimexpr\ht\pandoc@box+\dp\pandoc@box\relax}%
  \Gscale@div\@tempb{\linewidth}{\wd\pandoc@box}%
  \ifdim\@tempb\p@<\@tempa\p@\let\@tempa\@tempb\fi% select the smaller of both
  \ifdim\@tempa\p@<\p@\scalebox{\@tempa}{\usebox\pandoc@box}%
  \else\usebox{\pandoc@box}%
  \fi%
}
% Set default figure placement to htbp
\def\fps@figure{htbp}
\makeatother
\setlength{\emergencystretch}{3em} % prevent overfull lines
\providecommand{\tightlist}{%
  \setlength{\itemsep}{0pt}\setlength{\parskip}{0pt}}
\usepackage{booktabs}
\usepackage{float}
\usepackage{caption}
\usepackage{fancyhdr}
\usepackage{graphicx}
\usepackage{hyperref}
\usepackage{longtable}
\usepackage{array}
\hypersetup{colorlinks=true, linkcolor=blue, urlcolor=blue}
\pagestyle{fancy}
\fancyhead[L]{Education au Nigeria -- GHS Wave 4}
\fancyhead[R]{\thepage}
\fancyfoot[C]{}
\captionsetup{font=small, labelfont=bf, justification=centering}
\usepackage{booktabs}
\usepackage{longtable}
\usepackage{array}
\usepackage{multirow}
\usepackage{wrapfig}
\usepackage{float}
\usepackage{colortbl}
\usepackage{pdflscape}
\usepackage{tabu}
\usepackage{threeparttable}
\usepackage{threeparttablex}
\usepackage[normalem]{ulem}
\usepackage{makecell}
\usepackage{xcolor}
\usepackage{bookmark}
\IfFileExists{xurl.sty}{\usepackage{xurl}}{} % add URL line breaks if available
\urlstyle{same}
\hypersetup{
  pdftitle={Éducation et Scolarisation au Nigeria},
  pdfauthor={Analyse réalisée sous R},
  hidelinks,
  pdfcreator={LaTeX via pandoc}}

\title{Éducation et Scolarisation au Nigeria}
\usepackage{etoolbox}
\makeatletter
\providecommand{\subtitle}[1]{% add subtitle to \maketitle
  \apptocmd{\@title}{\par {\large #1 \par}}{}{}
}
\makeatother
\subtitle{Analyse de l'Enquête Générale sur les Ménages (GHS) --- Wave
4}
\author{Analyse réalisée sous R}
\date{17 mars 2026}

\begin{document}
\maketitle

{
\setcounter{tocdepth}{3}
\tableofcontents
}
\section{Introduction}\label{introduction}

Le capital humain constitue l'un des leviers les plus déterminants du
développement économique et social. Au Nigeria, pays le plus peuplé
d'Afrique subsaharienne avec plus de 200 millions d'habitants, les
inégalités d'accès à l'éducation persistent à de nombreux niveaux :
entre les sexes, entre les zones de résidence (urbaine et rurale), et
entre les États de la fédération.

La présente analyse exploite les données de la \textbf{quatrième vague
de l'Enquête Générale sur les Ménages} (\emph{General Household Survey},
GHS Wave 4), collectée par le Bureau National des Statistiques du
Nigeria (NBS) en partenariat avec la Banque Mondiale. Cette enquête à
représentativité nationale offre des informations détaillées sur les
caractéristiques socio-démographiques des individus et des ménages,
notamment leur niveau d'instruction et leur statut scolaire.

L'objectif de ce rapport est de dresser un portrait statistique de la
situation éducative au Nigeria à partir de quatre axes d'analyse
complémentaires :

\begin{enumerate}
\def\labelenumi{\arabic{enumi}.}
\tightlist
\item
  La \textbf{distribution générale du niveau d'instruction} dans la
  population adulte ;
\item
  Les \textbf{disparités éducatives entre hommes et femmes} chez les
  adultes de 18 ans et plus ;
\item
  La \textbf{relation entre l'âge et le niveau d'instruction},
  indicateur des transformations éducatives intergénérationnelles ;
\item
  Les \textbf{écarts de scolarisation entre zones urbaines et rurales}
  pour les enfants de 6 à 17 ans ;
\item
  La \textbf{cartographie de l'analphabétisme} à travers les 36 États et
  le Territoire Fédéral de la Capitale (FCT).
\end{enumerate}

Chaque section présente les visualisations graphiques correspondantes,
accompagnées des tests statistiques appropriés permettant d'évaluer la
significativité des écarts observés.

\section{Données et Méthodes}\label{donnuxe9es-et-muxe9thodes}

\subsection{Source des données}\label{source-des-donnuxe9es}

Les données utilisées proviennent de la \textbf{GHS Wave 4}
(\emph{General Household Survey, Panel --- Wave 4}), enquête
transversale représentative conduite au Nigeria. Elles couvrent les
ménages ruraux et urbains répartis sur l'ensemble des États du pays.
Trois fichiers principaux ont été mobilisés :

\begin{itemize}
\tightlist
\item
  \textbf{\texttt{secta\_harvestw4.dta}} : informations au niveau ménage
  (identifiant, zone de résidence, État) ;
\item
  \textbf{\texttt{sect1\_harvestw4.dta}} : caractéristiques
  démographiques individuelles (sexe, âge, etc.) ;
\item
  \textbf{\texttt{sect2\_harwestw4.dta}} : informations relatives à
  l'éducation et à la scolarisation.
\end{itemize}

\subsection{Traitement et nettoyage des
données}\label{traitement-et-nettoyage-des-donnuxe9es}

Les variables d'intérêt ont été recodées et structurées comme suit :

\begin{itemize}
\tightlist
\item
  \textbf{Niveau d'instruction} (\texttt{niveau\_educ}) : variable
  catégorielle ordonnée regroupant les modalités \emph{Aucun},
  \emph{Primaire}, \emph{Secondaire}, \emph{Supérieur} ;
\item
  \textbf{Sexe} (\texttt{s1q2}) : variable binaire recodée en
  \emph{Homme} / \emph{Femme} après décodage des étiquettes Stata via
  \texttt{haven::zap\_labels()} ;
\item
  \textbf{Scolarisation} (\texttt{s2aq6}) : variable dichotomique
  \emph{Oui} / \emph{Non} indiquant si l'enfant est actuellement inscrit
  à l'école ;
\item
  \textbf{Zone de résidence} (\texttt{sector}) : variable binaire
  \emph{Urbain} / \emph{Rural} issue du fichier ménage.
\end{itemize}

Les analyses portent sur les adultes de \textbf{18 ans et plus} pour les
questions relatives au niveau d'instruction, et sur les \textbf{enfants
de 6 à 17 ans} pour les analyses de scolarisation.

\subsection{Méthodes statistiques}\label{muxe9thodes-statistiques}

Les méthodes statistiques mobilisées sont les suivantes :

\begin{itemize}
\tightlist
\item
  \textbf{Statistiques descriptives} : distributions de fréquences,
  proportions, intervalles de confiance à 95 \% ;
\item
  \textbf{Test du Khi-deux} d'indépendance pour les relations entre
  variables catégorielles ;
\item
  \textbf{V de Cramér} pour mesurer l'intensité de l'association entre
  deux variables qualitatives ;
\item
  \textbf{Test de Kruskal-Wallis} (alternative non-paramétrique à
  l'ANOVA) pour comparer la distribution de l'âge entre les groupes
  éducatifs ;
\item
  \textbf{Visualisations} : barplots, boxplots, graphiques à barres
  empilées 100 \%, graphique à barres avec intervalles de confiance,
  heatmap par quintile.
\end{itemize}

Tous les tests sont réalisés au seuil de signification de \textbf{5 \%
(α = 0,05)}.

\section{Distribution du Niveau
d'Instruction}\label{distribution-du-niveau-dinstruction}

\subsection{Présentation graphique}\label{pruxe9sentation-graphique}

\begin{figure}[H]

{\centering \includegraphics[width=0.9\linewidth]{Rapport_final_files/figure-latex/task8-barplot-1} 

}

\caption{Distribution du niveau d'instruction dans la population enquêtée (GHS Nigeria, Wave 4). Les barres représentent le nombre d'individus pour chaque catégorie éducative. La palette de couleurs Mako (viridis) permet une lecture aisée des modalités.}\label{fig:task8-barplot}
\end{figure}

\subsection{Commentaire}\label{commentaire}

La \textbf{Figure 1} présente la répartition des individus de
l'échantillon selon leur niveau d'instruction. Plusieurs enseignements
peuvent être tirés de cette distribution :

\begin{itemize}
\tightlist
\item
  La catégorie \textbf{``Aucun''} (sans instruction) représente la
  modalité la plus fréquente, témoignant d'une proportion importante
  d'adultes n'ayant jamais été scolarisés. Ce résultat est cohérent avec
  les statistiques nationales publiées par l'UNESCO qui place le taux
  d'analphabétisme des adultes au Nigeria parmi les plus élevés
  d'Afrique subsaharienne.
\item
  Le \textbf{niveau primaire} constitue le deuxième groupe le plus
  représenté, indiquant que pour une partie significative de la
  population, la scolarisation s'est arrêtée avant l'achèvement du cycle
  primaire.
\item
  Les niveaux \textbf{secondaire} et \textbf{supérieur} restent
  minoritaires, reflétant les goulots d'étranglement qui caractérisent
  le système éducatif nigérian, particulièrement marqués dans les zones
  rurales et dans certaines régions du nord du pays.
\end{itemize}

Cette distribution en forme de pyramide inversée soulève des
préoccupations importantes quant à la disponibilité d'une main-d'œuvre
qualifiée, facteur essentiel pour la transformation économique du
Nigeria.

\section{Disparités Éducatives selon le
Sexe}\label{disparituxe9s-uxe9ducatives-selon-le-sexe}

\subsection{Distribution comparée}\label{distribution-comparuxe9e}

\begin{figure}[H]

{\centering \includegraphics[width=0.9\linewidth]{Rapport_final_files/figure-latex/task9-sexe-1} 

}

\caption{Répartition du niveau d'instruction par sexe chez les adultes de 18 ans et plus (GHS Nigeria, Wave 4). Graphique à barres empilées 100 \%. Chaque barre représente la composition éducative de chaque groupe de sexe, permettant une comparaison directe des proportions.}\label{fig:task9-sexe}
\end{figure}

\subsection{Résultats statistiques}\label{ruxe9sultats-statistiques}

\begingroup\fontsize{10}{12}\selectfont

\begin{longtable}[t]{lc}
\caption{\label{tab:task9-stats}Résultats du test du Khi-deux et V de Cramér pour la relation entre sexe et niveau d'instruction (adultes 18 ans et plus)}\\
\toprule
Statistique & Valeur\\
\midrule
\cellcolor{gray!10}{Khi-deux (χ²)} & \cellcolor{gray!10}{81.944}\\
Degrés de liberté & 4\\
\cellcolor{gray!10}{p-valeur} & \cellcolor{gray!10}{<0.001}\\
V de Cramér & 0.0782\\
\bottomrule
\end{longtable}
\endgroup{}

\subsection{Commentaire}\label{commentaire-1}

La \textbf{Figure 2} illustre de manière frappante les inégalités
éducatives entre hommes et femmes au Nigeria. Les observations majeures
sont :

\begin{itemize}
\tightlist
\item
  Les \textbf{femmes} présentent une proportion nettement plus élevée
  dans la catégorie \emph{``Aucun''} (sans instruction), confirmant une
  discrimination éducative historique à l'égard des filles,
  particulièrement prononcée dans les régions du nord où des facteurs
  culturels, religieux et économiques ont longtemps freiné la
  scolarisation féminine.
\item
  Les \textbf{hommes} sont sur-représentés dans les niveaux secondaire
  et supérieur, reflétant des stratégies d'investissement scolaire
  différenciées au sein des ménages qui, en situation de contrainte
  budgétaire, ont tendance à privilégier la scolarisation des garçons.
\end{itemize}

Le \textbf{Tableau 1} confirme statistiquement ces observations : le
test du Khi-deux est hautement significatif (p \textless{} 0,001),
rejetant l'hypothèse d'indépendance entre sexe et niveau d'instruction.
Le \textbf{V de Cramér}, compris entre 0,1 et 0,3, indique une
association de \textbf{force modérée à modeste} --- la relation est
réelle mais d'autres facteurs (géographiques, socio-économiques) jouent
également un rôle important dans la détermination du niveau
d'instruction.

\section{Relation entre l'Âge et le Niveau
d'Instruction}\label{relation-entre-luxe2ge-et-le-niveau-dinstruction}

\subsection{Boxplot comparatif}\label{boxplot-comparatif}

\begin{figure}[H]

{\centering \includegraphics[width=0.9\linewidth]{Rapport_final_files/figure-latex/task10-boxplot-1} 

}

\caption{Distribution de l'âge selon le niveau d'instruction chez les adultes de 18 ans et plus (GHS Nigeria, Wave 4). La boîte représente l'interquartile (Q1–Q3), la ligne médiane indique la médiane, et les moustaches s'étendent jusqu'à 1,5 fois l'écart interquartile. Les points représentent les valeurs extrêmes.}\label{fig:task10-boxplot}
\end{figure}

\subsection{Résultats du test de
Kruskal-Wallis}\label{ruxe9sultats-du-test-de-kruskal-wallis}

\begingroup\fontsize{10}{12}\selectfont

\begin{longtable}[t]{llc}
\caption{\label{tab:task10-kruskal}Résultats du test de Kruskal-Wallis pour la comparaison de l'âge entre les groupes de niveau d'instruction}\\
\toprule
 & Statistique & Valeur\\
\midrule
\cellcolor{gray!10}{Kruskal-Wallis chi-squared} & \cellcolor{gray!10}{H de Kruskal-Wallis} & \cellcolor{gray!10}{1769.087}\\
df & Degrés de liberté & 4\\
\cellcolor{gray!10}{} & \cellcolor{gray!10}{p-valeur} & \cellcolor{gray!10}{<0.001}\\
\bottomrule
\end{longtable}
\endgroup{}

\subsection{Commentaire}\label{commentaire-2}

La \textbf{Figure 3} révèle une relation inverse entre l'âge médian et
le niveau d'instruction, signe d'une transformation éducative
intergénérationnelle au Nigeria :

\begin{itemize}
\tightlist
\item
  Les individus classés \textbf{``Aucun''} (sans instruction) présentent
  une médiane d'âge plus élevée, suggérant que le phénomène
  d'analphabétisme est plus fortement concentré dans les cohortes plus
  anciennes, nées dans une période où l'accès à l'école était très
  limité.
\item
  Les individus ayant atteint le niveau \textbf{supérieur} sont en
  moyenne plus jeunes, traduisant l'expansion progressive du système
  éducatif nigérian au cours des dernières décennies, notamment sous
  l'effet de politiques publiques telles que l'\emph{Universal Basic
  Education (UBE)} lancée en 1999.
\item
  La dispersion (amplitude des boîtes et longueur des moustaches) est
  particulièrement large pour la catégorie \emph{``Aucun''}, indiquant
  une hétérogénéité dans les profils des non-scolarisés.
\end{itemize}

Le \textbf{Tableau 2} confirme statistiquement ces différences : le test
de Kruskal-Wallis (retenu car la normalité des distributions d'âge par
groupe ne peut être supposée) est hautement significatif (p \textless{}
0,001). Ces résultats témoignent d'une \textbf{transition éducative} en
cours au Nigeria, encore incomplète mais réelle.

\section{Scolarisation des 6--17 ans : Urbain
vs.~Rural}\label{scolarisation-des-617-ans-urbain-vs.-rural}

\subsection{Taux de scolarisation avec intervalles de
confiance}\label{taux-de-scolarisation-avec-intervalles-de-confiance}

\begin{figure}[H]

{\centering \includegraphics[width=0.9\linewidth]{Rapport_final_files/figure-latex/task11-urb-rur-1} 

}

\caption{Taux de scolarisation des enfants de 6 à 17 ans selon la zone de résidence (GHS Nigeria, Wave 4). Les barres indiquent la proportion d'enfants actuellement scolarisés. Les intervalles de confiance à 95 \% sont représentés par les lignes d'erreur verticales.}\label{fig:task11-urb-rur}
\end{figure}

\subsection{Résultats statistiques}\label{ruxe9sultats-statistiques-1}

\begingroup\fontsize{10}{12}\selectfont

\begin{longtable}[t]{lc}
\caption{\label{tab:task11-stats}Résultats du test du Khi-deux et V de Cramér pour la relation entre zone de résidence et statut de scolarisation (enfants 6–17 ans)}\\
\toprule
Statistique & Valeur\\
\midrule
\cellcolor{gray!10}{Khi-deux (χ²)} & \cellcolor{gray!10}{280.566}\\
Degrés de liberté & 1\\
\cellcolor{gray!10}{p-valeur} & \cellcolor{gray!10}{<0.001}\\
V de Cramér & 0.1752\\
\bottomrule
\end{longtable}
\endgroup{}

\subsection{Commentaire}\label{commentaire-3}

La \textbf{Figure 4} met en évidence un \textbf{écart significatif de
scolarisation} entre les zones urbaines et rurales chez les enfants de 6
à 17 ans. Les intervalles de confiance à 95 \% ne se chevauchent pas, ce
qui constitue un premier indicateur visuel de la significativité
statistique de cet écart.

\begin{itemize}
\tightlist
\item
  Les enfants résidant en \textbf{zone urbaine} présentent un taux de
  scolarisation notablement plus élevé que leurs homologues ruraux.
  L'accès plus aisé aux établissements scolaires, la densité plus
  importante des infrastructures éducatives et les revenus moyens plus
  élevés des ménages urbains expliquent en grande partie cette
  supériorité.
\item
  En \textbf{zone rurale}, les freins à la scolarisation sont multiples
  : distances importantes vers les écoles, travaux agricoles impliquant
  les enfants, pauvreté des ménages et offre éducative insuffisante
  (enseignants peu nombreux, classes surpeuplées, manque de matériel
  pédagogique).
\end{itemize}

Ces résultats s'inscrivent dans la continuité des rapports de l'UNICEF
sur le Nigeria, qui identifient la fracture urbain-rural comme l'une des
principales dimensions des inégalités éducatives dans le pays.

Le \textbf{Tableau 3} confirme statistiquement l'association entre zone
de résidence et scolarisation : le test du Khi-deux est très
significatif (p \textless{} 0,001). La valeur du \textbf{V de Cramér}
indique la force de cette association, permettant de quantifier
l'ampleur de la disparité observée.

\section{Cartographie de l'Analphabétisme par
État}\label{cartographie-de-lanalphabuxe9tisme-par-uxe9tat}

\subsection{Heatmap des taux sans
instruction}\label{heatmap-des-taux-sans-instruction}

\begin{figure}[H]

{\centering \includegraphics[width=0.9\linewidth]{Rapport_final_files/figure-latex/task12-heatmap-1} 

}

\caption{Proportion d'adultes (18 ans et plus) sans instruction formelle par État du Nigeria (GHS Wave 4). Les États sont classés par taux croissant. La colorisation par quintile (jaune = quintile inférieur, rouge foncé = quintile supérieur) permet d'identifier visuellement les États les plus défavorisés en matière d'éducation.}\label{fig:task12-heatmap}
\end{figure}

\subsection{Commentaire}\label{commentaire-4}

La \textbf{Figure 5} constitue la visualisation la plus riche de ce
rapport en révélant la profonde \textbf{hétérogénéité géographique} de
l'analphabétisme au Nigeria. Plusieurs constats s'imposent :

\begin{itemize}
\item
  Les États du \textbf{quintile 5} (rouge foncé), concentrant les taux
  d'adultes sans instruction les plus élevés, se situent majoritairement
  dans la \textbf{zone géo-politique du Nord-Ouest et du Nord-Est} du
  Nigeria (Sokoto, Kebbi, Zamfara, Jigawa, etc.). Ces États partagent
  des caractéristiques communes : forte ruralité, traditions de mariage
  précoce freinant la scolarisation des filles, et contexte sécuritaire
  parfois défavorable à la fréquentation scolaire.
\item
  À l'opposé, les États du \textbf{quintile 1} (jaune), affichant les
  taux d'analphabétisme les plus faibles, correspondent généralement aux
  États du \textbf{Sud-Ouest et du Sud-Est} ainsi qu'aux grandes
  agglomérations (Lagos, Abuja FCT, Rivers). Ces territoires bénéficient
  d'une tradition de scolarisation plus ancienne, d'une urbanisation
  plus avancée et d'un développement économique plus marqué.
\item
  La \textbf{dispersion des taux} entre le quintile inférieur et le
  quintile supérieur traduit non pas un retard uniforme mais bien deux
  Nigerias éducatifs dont les dynamiques et les défis divergent
  profondément.
\end{itemize}

Ces résultats appellent à des politiques publiques ciblées et
différenciées plutôt qu'une approche uniforme : si les États du Sud
nécessitent des actions de consolidation et d'amélioration qualitative,
ceux du Nord requièrent des interventions structurelles massives en
matière d'offre scolaire, de lutte contre la pauvreté et de changement
des normes sociales.

\section{Synthèse et Conclusion}\label{synthuxe8se-et-conclusion}

\subsection{Bilan des principaux
résultats}\label{bilan-des-principaux-ruxe9sultats}

\begingroup\fontsize{9}{11}\selectfont

\begin{longtable}[t]{>{\raggedright\arraybackslash}p{3cm}>{\raggedright\arraybackslash}p{6cm}>{\raggedright\arraybackslash}p{6cm}}
\caption{\label{tab:synthese-tab}Tableau de synthèse des principaux résultats et implications de politique publique}\\
\toprule
Axe d'analyse & Résultat principal & Implication de politique\\
\midrule
\cellcolor{gray!10}{Distribution générale} & \cellcolor{gray!10}{Forte prévalence du niveau 'Aucun' ; minorité ayant atteint l'enseignement supérieur} & \cellcolor{gray!10}{Nécessité d'alphabétisation des adultes et de renforcement du primaire}\\
Genre (adultes 18+) & Inégalité H/F significative (χ² p < 0,001) ; V de Cramér modéré & Politiques de scolarisation ciblant les filles, notamment dans le Nord\\
\cellcolor{gray!10}{Âge vs. éducation} & \cellcolor{gray!10}{Relation inverse âge-éducation (H Kruskal p < 0,001) : transition éducative en cours} & \cellcolor{gray!10}{Continuité des réformes : effets positifs visibles dans les jeunes cohortes}\\
Scolarisation 6-17 ans & Écart urbain-rural significatif (χ² p < 0,001) ; meilleur accès en zone urbaine & Investissement en infrastructures scolaires rurales prioritaire\\
\cellcolor{gray!10}{Géographie de l'analphabétisme} & \cellcolor{gray!10}{Gradient Nord-Sud marqué ; concentration de l'analphabétisme dans le Nord-Ouest} & \cellcolor{gray!10}{Interventions différenciées par État ; urgence dans les États du Nord}\\
\bottomrule
\end{longtable}
\endgroup{}

\subsection{Conclusion générale}\label{conclusion-guxe9nuxe9rale}

Cette analyse des données GHS Wave 4 confirme que la situation éducative
au Nigeria, bien qu'en amélioration au fil des générations, reste
marquée par des inégalités profondes selon le sexe, le milieu de
résidence et la localisation géographique.

Les résultats obtenus s'inscrivent dans une perspective de
\textbf{transition éducative} : les cohortes les plus jeunes sont en
moyenne mieux éduquées que leurs aînées, traduisant les effets cumulés
des politiques d'éducation universelle mises en œuvre depuis les années
2000. Toutefois, cette progression est loin d'être homogène, et les
disparités structurelles --- notamment entre le Nord et le Sud du pays
--- demeurent préoccupantes.

Pour atteindre les \textbf{Objectifs de Développement Durable (ODD 4)}
relatifs à l'éducation de qualité, le Nigeria devra intensifier ses
efforts sur plusieurs fronts : augmenter le budget de l'éducation
(actuellement inférieur aux recommandations internationales de 20

\end{document}
